\section{财富水平影响赚取–消费决策}\label{chap3_1:results_BHV}
猕猴在“赚取–消费”范式中的选择同时取决于选项条件和已累积的财富：在更高的赚取速率、更低的奖励量、以及更低的累积财富条件下，猕猴有更高概率选择赚取选项（图~\ref{fig1_3:p_earn}）。这一结果表明猕猴不仅习得了不同选项的价值，还能根据当前的财富水平调整赚取–消费偏好。

\begin{figure}[H]
    \centering
    \includegraphics[width=\textwidth]{Figure1_3_P_Earn}
    \bicaption{\enspace 赚取比例随财富水平的变化}{\enspace Earning choice proportion as a function of wealth}
    \label{fig1_3:p_earn}
\end{figure}
\fignotetext{图中展示了12种选项条件下，选择赚取的试次比例随财富水平（横轴，代币数1-24）的变化。颜色编码选项条件的客观价值比（与图~\ref{fig2_2:option_matrix}中的颜色梯度一致），偏红表示赚取价值相对更高的条件，偏绿表示消费价值相对更高的条件；亮度编码财富水平，颜色越深表示财富水平越低，颜色越浅表示财富水平越高。}{The proportion of earning choices across 12 option conditions is shown as a function of wealth level (x-axis, token count 1-24). Color encodes the objective value ratio of each condition (consistent with the color gradient in Figure~\ref{fig2_2:option_matrix}): redder colors indicate conditions in which earning value is relatively higher, and greener colors indicate conditions in which spending value is relatively higher. Brightness encodes wealth level, with darker colors indicating lower wealth and lighter colors indicating higher wealth.}

为在不预设函数形式的前提下量化财富对选项价值的调节作用，我们采用非参数方法（详见~\ref{chap2:methods_BHV_nonparametric}~节），估计了7种选项（3种赚取速率与4种消费奖励量）在24个财富水平上的主观权重。结果显示，赚取选项（3条红色曲线）的权重随财富增加而系统性下降，消费选项（4条绿色曲线）的权重则随财富增加而上升（图~\ref{fig1_4:subjective_weights}）。这一结果符合边际效用递减原理：随着财富增长，赚取额外代币的主观价值降低，而消费代币以获取即时奖励的吸引力上升。

\begin{figure}[H]
    \centering
    \includegraphics[width=\textwidth]{Figure1_4_SubjectiveWeights}
    \bicaption{\enspace 主观权重}{\enspace Subjective weights}
    \label{fig1_4:subjective_weights}
\end{figure}
\fignotetext{每张图展示7个选项的主观权重（$V^{\mathrm{Non}}$）随财富水平（横轴，代币数1–24）的变化。3条红色曲线代表赚取选项，颜色越深表示赚取速率越高；4条绿色曲线代表消费选项，颜色越深表示奖励量越大。曲线经平滑正则化处理（详见~\ref{chap2:methods_BHV_nonparametric}~节）。}{Each panel shows subjective weights of seven options ($V^{\mathrm{Non}}$) as a function of wealth (x-axis, token count 1-24). Three red curves represent earning options, with darker colors indicating higher earning rates; four green curves represent spending options, with darker colors indicating larger reward sizes. Curves were regularized across adjacent wealth levels (see Section~\ref{chap2:methods_BHV_nonparametric}).}

为刻画上述选择模式背后的计算机制，我们构建了基于边际效用理论的参数模型（~\ref{chap2:methods_BHV_utility}~节），用包括边际效用曲率指数 $\alpha$ 在内的四个自由参数解释猕猴的选择模式。模型估计出的主观价值之差 $\Delta V = V_E - V_S$ 能解释猕猴的选择偏好（图~\ref{fig1_5:psychometric}）。用留一交叉验证比较自由 $\alpha$ 模型（允许边际效用随财富递减）与固定 $\alpha$ 模型（假设线性边际效用，$\alpha = 1$）对猕猴行为的解释力，结果表明，自由 $\alpha$ 模型在三只猕猴中均显著优于固定 $\alpha$ 模型（图~\ref{fig1_6:model_comparison}）。进一步对每个实验日数据单独拟合效用模型，三只猕猴的 $\alpha$ 呈现集中分布模式，中位数分别为0.73、0.81和0.79，均显著低于1（图~\ref{fig1_7:alpha_dist}），验证了效用模型的稳健性。仅含4个自由参数的效用模型所导出的主观价值与选择模式，与拥有168个自由参数的非参数模型高度一致（图~\ref{fig1_8:model_pred}），表明低维参数模型已充分捕获了财富依赖的经济决策的核心结构。

\begin{figure}[H]
    \centering
    \includegraphics[width=\textwidth]{Figure1_5_PsychometricCurve}
    \bicaption{\enspace 心理物理曲线}{\enspace Psychometric curves}
    \label{fig1_5:psychometric}
\end{figure}
\fignotetext{赚取比例随效用模型估计出的主观价值差 $\Delta V = V_E - V_S$ 的变化（简洁起见，此处省略常数偏移项 $b$ 与选择敏感性参数 $\gamma$，但两者均已纳入实际计算）。每个数据点代表一种"选项条件$\times$财富"组合下的赚取比例，点的大小与该组合的试次数成正比。颜色沿用图~\ref{fig1_3:p_earn}的选项条件配色，偏红和偏绿分别表示客观上赚取价值更高和消费价值更高的条件，亮度编码当前财富水平（低财富水平为深色，高财富水平为浅色）。虚线为模型参数导出的心理物理曲线。}{Proportion of earn choices as a function of the subjective value difference $\Delta V = V_E - V_S$ estimated by the utility model (the bias term $b$ and choice sensitivity $\gamma$ were omitted from the formula for clarity but included in the computation). Each data point represents the earn proportion for one option-condition $\times$ wealth combination, with point size proportional to trial count. Colors follow the convention in Figure~\ref{fig1_3:p_earn}: redder colors indicate conditions in which earning value is higher, and greener colors indicate conditions in which spending value is higher; brightness encodes wealth level, with darker shades indicating lower wealth. The dashed curve is the psychometric curve derived from the utility model parameters.}

\begin{figure}[H]
    \centering
    \includegraphics[width=0.5\textwidth]{Figure1_6_Model_Comparison}
    \bicaption{\enspace 行为模型比较}{\enspace Behavioral model comparison}
    \label{fig1_6:model_comparison}
\end{figure}
\fignotetext{采用留一交叉验证评估模型泛化能力。以单个实验日为留存单元，用其余所有实验日的数据训练模型，在留存实验日数据上计算预测损失（详见\ref{chap2:methods_BHV_LOO}节）。纵轴为两模型的平均损失之差（$\Delta L = L^{\mathrm{fixed}}_{\mathrm{test}} - L^{\mathrm{free}}_{\mathrm{test}}$），正值表示自由 $\alpha$ 模型的样本外预测优于固定 $\alpha$ 模型。小提琴图展示各实验的 $\Delta L$ 分布，中线和端点分别对应中位数与四分位距。三只猕猴的自由 $\alpha$ 模型均显著优于固定 $\alpha$ 模型（猕猴W：52个实验日，$p = 4.2 \times 10^{-10}$；猕猴U：63个实验日，$p = 7.2 \times 10^{-12}$；猕猴T：64个实验日，$p = 3.5 \times 10^{-12}$；Wilcoxon符号秩检验）。}{Leave-one-session-out cross-validation was used to assess model generalization. On each fold, models were trained on all but one session and evaluated on the held-out session (see Section~\ref{chap2:methods_BHV_LOO}). The y-axis shows the per-trial average loss difference ($\Delta L = L_{\mathrm{fixed}} - L_{\mathrm{free}}$); positive values indicate better out-of-sample prediction by the free-$\alpha$ model. Violin plots show the distribution of $\Delta L$ across sessions, with lines indicating the median and interquartile range. The free-$\alpha$ model significantly outperformed the fixed-$\alpha$ model in all three monkeys (Monkey W: 52 sessions, $p = 4.2 \times 10^{-10}$; Monkey U: 63 sessions, $p = 7.2 \times 10^{-12}$; Monkey T: 64 sessions, $p = 3.5 \times 10^{-12}$; Wilcoxon signed-rank test).}

\begin{figure}[H]
    \centering
    \includegraphics[width=\textwidth]{Figure1_7_Alpha_Distribution}
    \bicaption{\enspace $\alpha$ 分布}{\enspace Distribution of fitted $\alpha$ across sessions}
    \label{fig1_7:alpha_dist}
\end{figure}
\fignotetext{每个实验日独立拟合自由 $\alpha$ 模型，图中展示三只猕猴所有实验日的 $\alpha$ 估计值分布（组距0.05）。蓝色虚线标注各猕猴$\alpha$分布的中位数：猕猴W为0.73，猕猴U为0.81，猕猴T为0.79。三只猕猴的 $\alpha$ 中位数均显著低于1（Wilcoxon符号秩检验，零假设 $\alpha = 1$，$p \approx 0$，低于数值精度）。}{The free-$\alpha$ model was fitted independently to each session, and the distribution of estimated $\alpha$ across all sessions is shown for each monkey (bin width = 0.05). Blue dashed lines indicate the median $\alpha$ for each monkey: Monkey W (52 sessions), 0.73; Monkey U (63 sessions), 0.81; Monkey T (64 sessions), 0.79. The median $\alpha$ was significantly below 1 in all three monkeys (Wilcoxon signed-rank test, null hypothesis $\alpha = 1$, $p \approx 0$ below numerical precision).}

\begin{figure}[H]
    \centering
    \includegraphics[width=\textwidth]{Figure1_8_BHV_Prediction}
    \bicaption{\enspace 效用模型预测的选择模式与主观价值}{\enspace Choice patterns and subjective values predicted by the utility model}
    \label{fig1_8:model_pred}
\end{figure}
\fignotetext{\textbf{(A) 观测与预测的选择模式}。赚取比例随财富水平的变化，颜色与亮度沿用图~\ref{fig1_3:p_earn}的约定。左图：实际赚取比例；右图：参数模型预测的赚取比例。\\ \textbf{(B) 观测与预测的主观价值差}。$\Delta V$ 随财富水平的变化。左图：非参数模型估计的 $\Delta V$（168个自由参数），作为高维参考基准；右图：参数模型导出的 $\Delta V$（4个自由参数）。}{\textbf{(A) Observed and predicted choice patterns}. Earning proportion as a function of wealth level; colors and brightness follow the convention in Figure~\ref{fig1_3:p_earn}. Left: observed earning proportion. Right: earning proportion predicted by the parametric model. \\\textbf{(B) Observed and predicted value difference}. $\Delta V$ as a function of wealth level. Left: $\Delta V$ estimated by the non-parametric model (168 free parameters), serving as the high-dimensional reference. Right: $\Delta V$ derived from the parametric model (4 free parameters).}

此外，效用模型估计的选项价值差的绝对值（ $|\Delta V|$ ）还能预测反应时间：在 $|\Delta V|$ 较小的试次中，赚取与消费选项的价值接近，反应时间也显著更长（图~\ref{fig1_9:rt_dist}）。

\begin{figure}[H]
    \centering
    \includegraphics[width=\textwidth]{Figure1_9_RT}
    \bicaption{\enspace 反应时分布}{\enspace Distribution of reaction time}
    \label{fig1_9:rt_dist}
\end{figure}
\fignotetext{试次按 $|\Delta V|$ 的中位数分为"简单"（$|\Delta V|$ 高于中位数，橙色）和"困难"（$|\Delta V|$ 低于中位数，蓝色）两组。直方图展示首次选择的反应时间分布，同色竖线标示各组中位数。价值差异较小的试次（蓝色，"困难"组）中反应时间更长（Wilcoxon秩和检验，猕猴W：$p = 1.3 \times 10^{-191}$；猕猴U：$p \approx 0$，低于数值精度；猕猴T：$p = 4.8 \times 10^{-223}$）。}{Trials were split at the median $|\Delta V|$ into "easy" ($|\Delta V|$ above median, orange) and "hard" ($|\Delta V|$ below median, blue) groups. Histograms show the reaction time distribution for first choices, with vertical lines marking the median of each group. Trials with smaller value differences (blue, "hard" group) elicited significantly longer decision times (Wilcoxon rank-sum test, Monkey W: $p = 1.3 \times 10^{-191}$; Monkey U: $p \approx 0$ below numerical precision; Monkey T: $p = 4.8 \times 10^{-223}$).}

这些结果共同表明，猕猴的赚取–消费选择同时依赖于选项价值与当前财富状态，效用模型能稳定解释猕猴的选择偏好和决策时间，说明猕猴的行为模式可以在边际效用递减框架内得到解释。
